% BHU PYQ -- Transcribed & Corrected
% Paper: BHUVA-13: Ayurveda -- Traditional Health Care & Management Principles
% Exam: U.G. Semester I Examination, 2024-25 (NEP)
% Category: Value Added Course (VAC)

\documentclass[11pt,a4paper]{article}
\usepackage[a4paper,top=2.5cm,bottom=2.5cm,left=2.8cm,right=2.8cm,headheight=14pt]{geometry}
\usepackage{amsmath,amssymb}
\usepackage[T1]{fontenc}
\usepackage[utf8]{inputenc}
\usepackage{lmodern,microtype,fancyhdr,enumitem,hyperref,lastpage,parskip}

\hypersetup{
    colorlinks=true,
    urlcolor=black,
    linkcolor=black,
    pdftitle={BHUVA13: Ayurveda -- BHU Sem I 2024-25 (NEP VAC)}
}

\pagestyle{fancy}
\fancyhf{}
\renewcommand{\headrulewidth}{0.4pt}
\renewcommand{\footrulewidth}{0.4pt}
\fancyhead[L]{\small\textit{Banaras Hindu University}}
\fancyhead[R]{\small\textit{BHUVA-13: Ayurveda (VAC)}}
\fancyfoot[C]{\small Page \thepage\ of \pageref{LastPage}}

\newlist{parts}{enumerate}{2}
\setlist[parts,1]{label=(\Alph*),leftmargin=*,itemsep=3pt,topsep=2pt}
\newcommand{\pts}[1]{\hfill\textit{[#1]}}

\begin{document}

\begin{center}
  {\large\bfseries BANARAS HINDU UNIVERSITY}\\[2pt]
  {\normalsize Under Graduate --- Semester I Examination, 2024-25 (NEP)}\\[6pt]
  \rule{0.8\linewidth}{0.5pt}\\[6pt]
  {\large\bfseries Value Added Course (VAC)}\\[4pt]
  {\large\bfseries Paper Code: BHUVA-13 : Ayurveda --- Traditional Health Care \& Management Principles}\\[4pt]
  {\normalsize Time: 1.5 Hours\hfill Maximum Marks: 70}
\end{center}

\rule{\linewidth}{0.4pt}

\smallskip
\noindent\textit{Instructions: Answer all 70 Multiple Choice Questions. Each question carries equal marks (1 mark each).}

\bigskip

\noindent\textbf{Question 1.} Charvaka Darshana comes under:\pts{1}
\begin{parts}
  \item Aastika Darshan
  \item Nastika Darshan
  \item Aastika-Nastika Darshan
  \item None of the above
\end{parts}

\medskip

\noindent\textbf{Question 2.} Which one is NOT included in Shad Padartha of Vaisheshik Darshan?\pts{1}
\begin{parts}
  \item Guna
  \item Dravya
  \item Karma
  \item Pada
\end{parts}

\medskip

\noindent\textbf{Question 3.} Yoga Sutras were expounded by:\pts{1}
\begin{parts}
  \item Maharshi Charaka
  \item Maharshi Patanjali
  \item Maharshi Sushruta
  \item Maharshi Kapila
\end{parts}

\medskip

\noindent\textbf{Question 4.} In Maheshwar Sutra, 'Ik' is a:\pts{1}
\begin{parts}
  \item Pratyahara
  \item Sandhi
  \item Samasa
  \item Dhatu
\end{parts}

\medskip

\noindent\textbf{Question 5.} Vedanta Darshan comes under which of the following:\pts{1}
\begin{parts}
  \item Aastika Darshan
  \item Nastika Darshan
  \item Both Aastika and Nastika
  \item Neither
\end{parts}

\medskip

\noindent\textbf{Question 6.} The term Darshana is derived from:\pts{1}
\begin{parts}
  \item Drish Dhatu
  \item Vid Dhatu
  \item Path Dhatu
  \item Kri Dhatu
\end{parts}

\medskip

\noindent\textbf{Question 7.} Ashtangahrudaya is written by:\pts{1}
\begin{parts}
  \item Charaka
  \item Sushruta
  \item Vagbhata
  \item Bhavamishra
\end{parts}

\medskip

\noindent\textbf{Question 8.} Who are the Devabhishak (Physicians of Gods)?\pts{1}
\begin{parts}
  \item Brahma
  \item Vishnu
  \item Ashwini Kumaras
  \item Shiva
\end{parts}

\medskip

\noindent\textbf{Question 9.} Total number of Sharirik Doshas in Ayurveda are:\pts{1}
\begin{parts}
  \item 2
  \item 3 (Vata, Pitta, Kapha)
  \item 4
  \item 5
\end{parts}

\medskip

\noindent\textbf{Question 10.} Total number of Manasik Doshas are:\pts{1}
\begin{parts}
  \item 2 (Raja, Tama)
  \item 3
  \item 4
  \item 5
\end{parts}

\medskip

\noindent\textbf{Question 11.} Vata, Pitta, and Kapha are together known as:\pts{1}
\begin{parts}
  \item Tridosha
  \item Triguna
  \item Trimala
  \item Trikala
\end{parts}

\medskip

\noindent\textbf{Question 12.} Which of the following compendia belongs to Brihattrayi?\pts{1}
\begin{parts}
  \item Charaka Samhita
  \item Sushruta Samhita
  \item Astanga Hridaya
  \item All of the above
\end{parts}

\medskip

\noindent\textbf{Question 13.} Which of the following compendia belongs to Laghuttrayi?\pts{1}
\begin{parts}
  \item Madhava Nidana
  \item Sharngadhara Samhita
  \item Bhava Prakasha
  \item All of the above
\end{parts}

\medskip

\noindent\textbf{Question 14.} Medicinal chemistry in Ayurveda primarily deals with:\pts{1}
\begin{parts}
  \item Dravyaguna and Rasa Shastra
  \item Anatomy
  \item Surgery
  \item Hygiene
\end{parts}

\medskip

\noindent\textbf{Question 15.} Yellowish discoloration of skin, eyes, urine, and faeces is characteristic of:\pts{1}
\begin{parts}
  \item Aggravated Pitta
  \item Aggravated Vata
  \item Aggravated Kapha
  \item Normal Rakta
\end{parts}

\medskip

\noindent\textbf{Question 16.} How many times should we brush our teeth a day according to Ayurveda?\pts{1}
\begin{parts}
  \item Once
  \item Twice (Dantdhavana)
  \item Thrice
  \item Four times
\end{parts}

\medskip

\noindent\textbf{Question 17.} ``Extra teeth and extra bones'' is a feature of:\pts{1}
\begin{parts}
  \item Aggravated Asthi Dhatu
  \item Diminished Shukra Dhatu
  \item Aggravated Medas Dhatu
  \item Diminished Rakta Dhatu
\end{parts}

\medskip

\noindent\textbf{Question 18.} Which among the following is/are included in the Shad Bhavas (six factors) which contribute to Garbha (embryo)?\pts{1}
\begin{parts}
  \item Matruja \& Pitruja
  \item Atmaja \& Rasaja
  \item Satttvaja \& Satmyaja
  \item All of the above
\end{parts}

\medskip

\noindent\textbf{Question 19.} Which of the following is a primary focus of Bhaishajya Kalpana in Ayurveda?\pts{1}
\begin{parts}
  \item Preparation, formulation, and application of medicinal substances
  \item Diagnostic tools for diseases
  \item Anatomy and physiology
  \item Dietary supplements only
\end{parts}

\medskip

\noindent\textbf{Question 20.} According to Ayurveda, health (Swasthya) is defined as a state of equilibrium of:\pts{1}
\begin{parts}
  \item Dosha, Agni, Dhatu, and Mala along with pleasant Mind and Senses
  \item Physical body only
  \item Absence of symptoms only
  \item High energy levels only
\end{parts}

\medskip

\noindent\textbf{Question 21.} Dinacharya in Ayurveda refers to:\pts{1}
\begin{parts}
  \item Daily healthy regimen
  \item Seasonal regimen
  \item Night regimen
  \item Dietary restrictions during illness
\end{parts}

\medskip

\noindent\textbf{Question 22.} Ritucharya refers to:\pts{1}
\begin{parts}
  \item Seasonal health regimen
  \item Daily routine
  \item Code of conduct
  \item Emergency management
\end{parts}

\medskip

\noindent\textbf{Question 23.} Sadvritta in Ayurveda promotes:\pts{1}
\begin{parts}
  \item Moral, ethical, and psychological health
  \item Surgical procedures
  \item Physical exercise only
  \item Fasting methods
\end{parts}

\medskip

\noindent\textbf{Question 24.} The main function of Agni in the human body is:\pts{1}
\begin{parts}
  \item Digestion, absorption, and metabolism
  \item Respiration
  \item Circulation of blood
  \item Excretion of sweat
\end{parts}

\medskip

\noindent\textbf{Question 25.} How many types of Agni are described in Ayurveda?\pts{1}
\begin{parts}
  \item 13 (1 Jatharagni, 5 Bhutagni, 7 Dhatvagni)
  \item 7
  \item 5
  \item 3
\end{parts}

\medskip

\noindent\textbf{Question 26.} Ojas is considered as the ultimate essence of:\pts{1}
\begin{parts}
  \item All seven Dhatus (Immunity and Vitality)
  \item Vata Dosha
  \item Blood only
  \item Waste products
\end{parts}

\medskip

\noindent\textbf{Question 27.} Prakriti (individual constitution) is determined at the time of:\pts{1}
\begin{parts}
  \item Conception
  \item Birth
  \item Childhood
  \item Adulthood
\end{parts}

\medskip

\noindent\textbf{Question 28.} Which Mahabhuta is dominant in Vata Dosha?\pts{1}
\begin{parts}
  \item Vayu (Air) and Akasha (Ether)
  \item Agni and Jala
  \item Prithvi and Jala
  \item Agni and Vayu
\end{parts}

\medskip

\noindent\textbf{Question 29.} Which Mahabhuta is dominant in Pitta Dosha?\pts{1}
\begin{parts}
  \item Agni (Fire) and Jala (Water)
  \item Vayu and Akasha
  \item Prithvi and Agni
  \item Jala and Earth
\end{parts}

\medskip

\noindent\textbf{Question 30.} Which Mahabhuta is dominant in Kapha Dosha?\pts{1}
\begin{parts}
  \item Prithvi (Earth) and Jala (Water)
  \item Vayu and Agni
  \item Akasha and Vayu
  \item Agni and Jala
\end{parts}

\medskip

\noindent\textbf{Question 31.} Panchakarma therapies are primarily used for:\pts{1}
\begin{parts}
  \item Bio-purification and detoxification of the body
  \item Surgical intervention
  \item Immediate pain relief
  \item Cosmetic enhancement
\end{parts}

\medskip

\noindent\textbf{Question 32.} Which of the following is NOT one of the classic Panchakarma procedures?\pts{1}
\begin{parts}
  \item Vamana (Emesis)
  \item Virechana (Purgation)
  \item Basti (Enema)
  \item Shirodhara (Oil pouring relaxation therapy)
\end{parts}

\medskip

\noindent\textbf{Question 33.} Nasya procedure involves administration of medicine through:\pts{1}
\begin{parts}
  \item Nasal route
  \item Oral route
  \item Rectal route
  \item Topical application
\end{parts}

\medskip

\noindent\textbf{Question 34.} Abhyanga refers to:\pts{1}
\begin{parts}
  \item Therapeutic full-body oil massage
  \item Herbal steam bath
  \item Medicated eye drops
  \item Fasting therapy
\end{parts}

\medskip

\noindent\textbf{Question 35.} Rasayana therapy in Ayurveda is aimed at:\pts{1}
\begin{parts}
  \item Rejuvenation, longevity, and enhancing memory/immunity
  \item Immediate cure of acute fever
  \item Surgical removal of tumors
  \item Purgation therapy
\end{parts}

\medskip

\noindent\textbf{Question 36.} Vajikarana therapy deals with:\pts{1}
\begin{parts}
  \item Reproductive health, virility, and aphrodisiacs
  \item Pediatric disorders
  \item Mental health disorders
  \item Eye diseases
\end{parts}

\medskip

\noindent\textbf{Question 37.} Which Rasa (taste) is dominant in Amla (sour) substances?\pts{1}
\begin{parts}
  \item Earth and Fire
  \item Water and Fire
  \item Air and Ether
  \item Earth and Water
\end{parts}

\medskip

\noindent\textbf{Question 38.} Total number of fundamental Rasas (tastes) recognized in Ayurveda is:\pts{1}
\begin{parts}
  \item 6 (Madhura, Amla, Lavana, Katu, Tikta, Kashaya)
  \item 4
  \item 5
  \item 7
\end{parts}

\medskip

\noindent\textbf{Question 39.} Madhura Rasa (Sweet taste) decreases which Dosha?\pts{1}
\begin{parts}
  \item Vata and Pitta
  \item Kapha
  \item Vata and Kapha
  \item Pitta and Kapha
\end{parts}

\medskip

\noindent\textbf{Question 40.} Katu Rasa (Pungent taste) increases which Doshas?\pts{1}
\begin{parts}
  \item Vata and Pitta
  \item Kapha and Pitta
  \item Vata and Kapha
  \item None
\end{parts}

\medskip

\noindent\textbf{Question 41.} What is the main cause of disease according to Ayurveda?\pts{1}
\begin{parts}
  \item Imbalance of Tridosha and Dhatus
  \item Lack of modern medicine
  \item Physical injury only
  \item Over-exercising only
\end{parts}

\medskip

\noindent\textbf{Question 42.} Ama in Ayurveda refers to:\pts{1}
\begin{parts}
  \item Toxic, undigested metabolic waste product
  \item Digested food nutrient
  \item Blood tissue
  \item Pure water
\end{parts}

\medskip

\noindent\textbf{Question 43.} Srotas in Ayurveda refers to:\pts{1}
\begin{parts}
  \item Micro and macro channels of circulation and transportation
  \item Bones
  \item Sensory organs
  \item Brain centers
\end{parts}

\medskip

\noindent\textbf{Question 44.} How many Dhatus (bodily tissues) are described in Ayurveda?\pts{1}
\begin{parts}
  \item 7 (Rasa, Rakta, Mamsa, Meda, Asthi, Majja, Shukra)
  \item 5
  \item 9
  \item 12
\end{parts}

\medskip

\noindent\textbf{Question 45.} Which Dhatu is responsible for nourishment of the body (Preenana)?\pts{1}
\begin{parts}
  \item Rasa Dhatu
  \item Rakta Dhatu
  \item Mamsa Dhatu
  \item Asthi Dhatu
\end{parts}

\medskip

\noindent\textbf{Question 46.} Which Dhatu provides structural support and stability to the body?\pts{1}
\begin{parts}
  \item Asthi Dhatu
  \item Meda Dhatu
  \item Rasa Dhatu
  \item Shukra Dhatu
\end{parts}

\medskip

\noindent\textbf{Question 47.} Ahara (Food), Nidra (Sleep), and Brahmacharya (Self-control) are known as:\pts{1}
\begin{parts}
  \item Trayopastambha (Three Pillars of Life)
  \item Tridosha
  \item Trimala
  \item Triguna
\end{parts}

\medskip

\noindent\textbf{Question 48.} Sattva, Rajas, and Tamas are known as:\pts{1}
\begin{parts}
  \item Trigunas (Mental attributes)
  \item Tridoshas
  \item Dhatus
  \item Rasas
\end{parts}

\medskip

\noindent\textbf{Question 49.} The father of Surgery in Ayurveda is considered to be:\pts{1}
\begin{parts}
  \item Maharshi Sushruta
  \item Maharshi Charaka
  \item Maharshi Vagbhata
  \item Maharshi Agnivesha
\end{parts}

\medskip

\noindent\textbf{Question 50.} The main text for Ayurvedic internal medicine (Kaya Chikitsa) is:\pts{1}
\begin{parts}
  \item Charaka Samhita
  \item Sushruta Samhita
  \item Astanga Hridaya
  \item Madhava Nidana
\end{parts}

\medskip

\noindent\textbf{Question 51.} The process of Swedana in Panchakarma refers to:\pts{1}
\begin{parts}
  \item Sudation / Herbal steam therapy
  \item Oil massage
  \item Bloodletting
  \item Nasal drops
\end{parts}

\medskip

\noindent\textbf{Question 52.} Raktamokshana is a procedure used for:\pts{1}
\begin{parts}
  \item Bloodletting therapy
  \item Emesis
  \item Purgation
  \item Enema
\end{parts}

\medskip

\noindent\textbf{Question 53.} Shirodhara is particularly beneficial for:\pts{1}
\begin{parts}
  \item Stress, anxiety, and insomnia
  \item Joint dislocations
  \item Skin burns
  \item Digestive disorders
\end{parts}

\medskip

\noindent\textbf{Question 54.} Triphala is a combination of three fruits, namely:\pts{1}
\begin{parts}
  \item Haritaki, Bibhitaki, and Amalaki
  \item Tulsi, Neem, and Ginger
  \item Pippali, Maricha, and Shunthi
  \item Ashwagandha, Shatavari, and Bala
\end{parts}

\medskip

\noindent\textbf{Question 55.} Trikatu consists of:\pts{1}
\begin{parts}
  \item Shunthi (Ginger), Maricha (Black Pepper), and Pippali (Long Pepper)
  \item Haritaki, Bibhitaki, and Amalaki
  \item Cardamom, Cinnamon, and Cloves
  \item Garlic, Onion, and Ginger
\end{parts}

\medskip

\noindent\textbf{Question 56.} Which herb is widely recognized in Ayurveda as an adaptogen and nervous system tonic?\pts{1}
\begin{parts}
  \item Ashwagandha
  \item Haridra
  \item Neem
  \item Senna
\end{parts}

\medskip

\noindent\textbf{Question 57.} Tulsi (\textit{Ocimum sanctum}) is primarily used in Ayurvedic remedies for:\pts{1}
\begin{parts}
  \item Respiratory disorders and immunity
  \item Fracture healing
  \item Laxative effect
  \item Sedative effect
\end{parts}

\medskip

\noindent\textbf{Question 58.} Which mineral/metal purification process is essential in Rasa Shastra?\pts{1}
\begin{parts}
  \item Shodhana
  \item Marana
  \item Amritikarana
  \item All of the above
\end{parts}

\medskip

\noindent\textbf{Question 59.} Bhasma in Ayurveda refers to:\pts{1}
\begin{parts}
  \item Calcined nano-sized metallic or mineral ash preparation
  \item Raw herb powder
  \item Liquid decoction
  \item Herbal ghee
\end{parts}

\medskip

\noindent\textbf{Question 60.} Asava and Arishta are examples of:\pts{1}
\begin{parts}
  \item Self-generated fermented biomedical formulations
  \item Medicated pastes
  \item Herbal teas
  \item Raw powders
\end{parts}

\medskip

\noindent\textbf{Question 61.} Kwatha in Ayurvedic pharmaceutics refers to:\pts{1}
\begin{parts}
  \item Herbal decoction made by boiling drugs in water
  \item Cold infusion
  \item Medicated oil
  \item Herbal pill
\end{parts}

\medskip

\noindent\textbf{Question 62.} Swarasa refers to:\pts{1}
\begin{parts}
  \item Fresh plant juice
  \item Dried herbal powder
  \item Distilled water
  \item Ferment solution
\end{parts}

\medskip

\noindent\textbf{Question 63.} Hima in Bhaishajya Kalpana refers to:\pts{1}
\begin{parts}
  \item Cold infusion of herbs
  \item Hot decoction
  \item Medicated jam
  \item Medicated butter
\end{parts}

\medskip

\noindent\textbf{Question 64.} Phanta refers to:\pts{1}
\begin{parts}
  \item Hot infusion prepared by pouring boiling water on herbs
  \item Cold juice
  \item Herbal paste
  \item Oil extract
\end{parts}

\medskip

\noindent\textbf{Question 65.} Kalka in Ayurvedic pharmacy means:\pts{1}
\begin{parts}
  \item Freshly ground herbal paste
  \item Herbal wine
  \item Mineral powder
  \item Clear liquid
\end{parts}

\medskip

\noindent\textbf{Question 66.} Churna refers to:\pts{1}
\begin{parts}
  \item Fine dry herbal powder
  \item Medicated oil
  \item Steam extract
  \item Fresh juice
\end{parts}

\medskip

\noindent\textbf{Question 67.} Taila and Ghrita preparations in Ayurveda belong to which category?\pts{1}
\begin{parts}
  \item Sneha Kalpana (Medicated Fats/Oils)
  \item Sandhana Kalpana
  \item Rasa Kalpana
  \item Churna Kalpana
\end{parts}

\medskip

\noindent\textbf{Question 68.} Lehya / Avaleha (e.g. Chyawanprash) is a type of:\pts{1}
\begin{parts}
  \item Semi-solid confectionary formulation
  \item Liquid extract
  \item Dry tablet
  \item Medicated ash
\end{parts}

\medskip

\noindent\textbf{Question 69.} Vati / Gutika in Ayurvedic pharmaceutics refers to:\pts{1}
\begin{parts}
  \item Medicated pills or tablets
  \item Herbal juices
  \item Medicated steam
  \item Fermented wines
\end{parts}

\medskip

\noindent\textbf{Question 70.} The goal of Swasthavritta in Ayurveda is:\pts{1}
\begin{parts}
  \item Maintenance of positive health and prevention of diseases
  \item Treatment of late-stage diseases
  \item Surgical amputations
  \item Commercial drug manufacturing
\end{parts}

\vfill
\rule{\linewidth}{0.4pt}
\begin{center}\small
\href{https://bhu-exam.vercel.app/}{Click here to visit our website}\\[4pt]
\href{https://me-aryan.vercel.app/}{Click here to visit our contributor}
\end{center}

\end{document}
